\lesson{3}{Sep 22 2021 Wen (11:05:57)}{Citizens in Action}{Unit 1}

\subsubsection*{What Are the Responsibilities of Citizenship}

Some of the responsibilities of citizens are required in law. We refer to these
as \bf{civic duties}. One civic duty is to obey laws. Everyone living in
the United States must abide by the laws, whether they are legal citizens or
not. However, citizens can legally act to change laws they feel are unjust. For
example, they could petition the government for a change in the law, contacting
their representatives and other citizens for support.

Responsibilities differ from duties because they are voluntary. The law does
not require citizens to complete these responsibilities, but people do have
certain expectations within their communities. Good citizens choose to meet
their community's expectations. Americans across the country have specific
ideas about what good citizenship means and what the expectations are.
Influencing government, such as contacting leaders or holding office, is one of
these expectations. People who are not legal U.S. citizens can also show good
citizenship by meeting many of these responsibilities and expectations.
However, remember that some of the responsibilities and duties, such as serving
on a jury and voting, are reserved only for legal citizens.

\subsubsection*{Why is Participation Important in the American Political System}

In American politics, is responsible for upholding our system of laws. Also,
justice is dependent on both the government and the people. The government:

\begin{itemize}
    \item Creates
    \item Enforces
    \item Settles
\end{itemize}

disputes over laws. It is up to the citizens to choose people carefully to
fulfill those roles in government, support the system of laws through obeying
them, and challenge unfair laws through legal means. Without this two-way
support of the system, representative democracy could dissolve into a system
that is not representative of the people, is abusive of the people, or becomes
an anarchy.

Tools such as:

\begin{itemize}
    \item EMail
    \item Blogs
    \item Podcasts
\end{itemize}

help connect people with the same goals and ideas across wide geographic areas.
People coordinate activities such as protests and demonstrations using these
tools.

Americans who are complacent to simply follow laws and pay taxes without
exercising responsible citizenship run the risk of allowing government to take
more of that responsibility for itself. People need to hold government
accountable for its actions in order to protect their interests in the laws.
While people have other responsibilities and most cannot devote all their time
to civic life, it is an essential responsibility for all citizens. Many
Americans develop new ideas for public policy, serve as watchdogs over
government, and inspire other citizens to meet their responsibilities.

A great challenge is that most citizens only focus on one or two issues, only
about half vote, and very few are active in political parties or interest
groups. Yet these latter groups raise millions of dollars to influence voters
and elected officials. The fewer people involved in directing that influence,
the less that government will represent the wishes of the general population.
To protect our representative democracy, more citizens need to speak up through
contacting officials, voting, and other forms of political participation.

\subsubsection*{How Have Models of Effective Citizenship Achieved Change?}

Citizenship is a lifelong responsibility to the common good of your political
community and the world. The common good depends on people taking action. Many
people in American history have acted on the responsibilities of citizenship to
achieve change. Individual action can inspire whole groups of people to get
involved, and people working together is a more powerful and effective way to
create change.

Some would say the most active citizens are those who devote their lives to the
military or to political office. Examples of great American presidents, other
elected officials, and military personnel abound in history. This timeline will
focus on a few citizens outside those areas. Each of these people was a leader
and part of a larger movement. They cannot claim sole credit for the effects of
their actions, but their responsible citizenship inspired many others and
helped create change.

\subsubsection*{How Have Models of Effective Citizenship Achieved Change}

Citizenship is a lifelong responsibility to the common good of your political
community and the world. The common good depends on people taking action. Many
people in American history have acted on the responsibilities of citizenship to
achieve change. Individual action can inspire whole groups of people to get
involved, and people working together is a more powerful and effective way to
create change.

\newpage
